% Appendix A — polished, paste-ready replacement for chapters/92-AppendixA.tex.
% TODO(Ben): bracket stiffness third diagonal — code uses chi2 (minimizeFlxPin.m:266-268),
% text eq (bktstiffness) states chi1. This file follows the CODE. Resolve before final.
% Cross-references (eq:*, tab:*, sec:*, ch:*) resolve inside the full document.

\chapter{Muscle Force and Joint Torque Calculations}\label{app:muscle_calcs}

This appendix documents the reference implementation of the force, torque, and
compliance equations used in Chapters~\ref{ch:methods}--\ref{ch:results}, as
they exist in the project's Matlab code, so that any number in the dissertation
can be traced to the code that produced it.

\section{BPA Force Model Implementation}

\subsection{Strain Measures}

The musculotendon path length $l_{LMT}$ is the sum of straight segments between
route via points, transformed into a common frame at each crossing. The
contractile strain subtracts the artificial tendon and both end-cap fittings
from the path length before normalization by the resting length,
\begin{equation}
    \varepsilon \;=\; \frac{l_{0} - \left(l_{LMT} - l_{ten} - 2\,l_{fitting} - \chi_{0}\right)}{l_{0}},
    \label{eq:app_strain}
\end{equation}
where $\chi_{0}$ is the constant length offset of
Section~\ref{sec:improved_model} (zero in the uncorrected model),
$l_{ten}$ the tendon length, and $l_{fitting} = \qty{25.4}{\mm}$ the end-cap
fitting length. The maximum free-load contraction at \qty{620}{\kPa} is
$\varepsilon_{620} = (l_{0}-l_{620})/l_{0}$, with $l_{620}$ the BPA length at
\qty{620}{\kPa} under no external load, and the relative strain is
$\varepsilon^{*} = \varepsilon / \varepsilon_{620}$.

\subsection{Normalized Force Surface}

The normalized isometric force surface of Equation~(\ref{eq:Fstar}) is
implemented in closed form,
\begin{equation}
    F^{*} = c_{0}\left(e^{-c_{1}\varepsilon^{*}} - 1\right)
            + P^{*}\,e^{-c_{2}(\varepsilon^{*})^{2}},
    \qquad F^{*}\le 0 \;\Rightarrow\; 0,
    \label{eq:app_fstar}
\end{equation}
with $P^{*} = P/\qty{620}{\kPa}$. The coefficients were fitted by nonlinear
least squares with least-absolute-residual robustness and are reported in
Table~\ref{tab:Fstar}. The code additionally holds a \qty{40}{\mm} fit that this
dissertation does not use ($c_0 = 0.1224$, $c_1 = 10.47$, $c_2 = 2.023$). A
spline interpolation of the manufacturer's datasheet curves was used as a
runtime alternative during model development, with force clamped to zero for
$\varepsilon^{*}>1$ or $F^{*}<0$; Equation~(\ref{eq:app_fstar}) superseded it.

\subsection{Maximum Force}

The resting-length-dependent maximum force of Equation~(\ref{eq:Fmax_3D}) uses
the fixed \qty{7.5}{\mm} end-contact offset. The fitted \qty{20}{\mm}
coefficients carry 95\% confidence intervals of
$a_1 \in [1.4745, 1.5011]$\,\si{\N\per\kPa} and
$a_2 \in [0.0238, 0.0258]$\,\si{\per\kPa\per\metre}. Total force is the product
$F = F^{*}(\varepsilon^{*},P^{*})\,F_{620}(l_{0})$, resolved into a vector along
the force direction.

\section{Muscle Path Length and Moment Arm}

Path length accumulates the Euclidean norm of each segment between via points,
applying the body-frame transformation whenever the path crosses between the
femur and tibia frames. The vector moment arm is implemented as the
perpendicular foot of the joint origin onto the line of action,
\begin{equation}
    \vec{r} \;=\; \vec{p} \;-\; \hat{u}\left(\hat{u}\cdot\vec{p}\right),
    \label{eq:app_perpfoot}
\end{equation}
where $\vec{p}$ runs from the joint origin to the muscle-path crossing point and
$\hat{u}$ is the unit force direction. This is the same quantity as the
cross-product projection of Equation~(\ref{eq:project_ma}); the planar scalar
moment arm about the joint axis is the magnitude of the in-plane components.
Torque about the joint $z$-axis is $\vec{M} = \vec{r}\times\vec{F}$, evaluated
only where the BPA is in tension; values in the compressive region
($\varepsilon < -0.02$) are discarded.

\section{Compliance Model Implementation}

The compliance-corrected muscle length entering Equation~(\ref{eq:Lm}) is
computed as $l_{m} = l_{LMT} - \chi_{0} - l_{ten} - 2\,l_{fitting}$. The bracket
stiffness matrix and the projected scalar compliance along the force direction
are evaluated per configuration as
$\mathbf{K}_{br} = \mathrm{diag}(\chi_{1},\chi_{2},\chi_{2})$ in the bracket
frame and $k_{\hat{u},br} = \left(\hat{u}^{\,T}\mathbf{K}_{br}^{-1}\hat{u}\right)^{-1}$.
The artificial tendon is modeled as an axial spring with stiffness
\begin{equation}
    k_{ten} \;=\; \frac{n\,A_{c}\,E}{L_{ten}},
    \label{eq:app_tendon}
\end{equation}
with $n=2$ parallel cable strands ($\times 6$ for the \qty{20}{\mm} assembly),
$A_{c} = \qty{1.51e-6}{\metre\squared}$ for the 19-strand cable, and
$E = \qty{193}{\giga\pascal}$. Series combination with the projected bracket
compliance gives $k_{eq} = \left(k_{\hat{u},br}^{-1} + k_{ten}^{-1}\right)^{-1}$,
as in Equation~(\ref{eq:equivalentstiffness}).

The scalar compliance displacement $\delta$ of
Equation~(\ref{eq:complianceforcebalance}) is found numerically with a
bracketed root solve on $[0,\, l_{m}-l_{620}]$,
\begin{equation}
    F_{620}(l_{0})\,F^{*}\!\left(\varepsilon^{*}(\delta),P^{*}\right) - k_{eq}\,\delta = 0,
    \label{eq:app_forcebalance}
\end{equation}
with $\delta \equiv 0$ when the uncorrected relative strain already exceeds
unity. The vector bracket deflection then follows from
$\boldsymbol{\epsilon}_{br} = \mathbf{K}_{br}^{-1}\left(\norm{\vec{F}}\hat{u}\right)$,
the cable stretch is $\norm{\vec{F}}/k_{ten}$, and the attachment point, path,
moment arm, and torque are recomputed. For the two-BPA flexor grouping, the
coupled balance shares one common force between both muscles and is solved
analogously on $[0, F_{620}]$.

\section{Wrapping-Loss Implementation}

For the pinned-knee extensor, Equation~(\ref{eq:deltaL}) is implemented
segment-wise over the wrap arc with radii $R_1 = \qty{12}{\mm}$ and
$R_2 = \qty{40}{\mm}$ and segment boundaries $\theta_{1} = \ang{27}$ and
$\theta_{2} = \ang{-23}$. For the biomimetic-knee \qty{20}{\mm} configuration,
the arc-length factor is the routed bend measure $b(\theta_{k})$, the arc length
swept by the wrap, giving
$\Delta l = \chi_{3}\, b(\theta_{k})\,(1-\varepsilon^{*})^{2}$, with a
hardcoded-geometry fallback for degenerate routes. The loss enters only the
force-producing strain; the kinematic strain used for prediction excludes it.
